% Generated by the blueprint tool (blueprint-pdf/1) from a blueprint snapshot; edit the
% Markdown sources, not this file.  Compile with `latexmk -pdf` or by running
% `pdflatex` twice (the second run settles the contents and the page references).
% Images, if any, are read from assets/ next to this file.
\documentclass[11pt]{article}
\usepackage[T1]{fontenc}
\usepackage[utf8]{inputenc}
\usepackage{lmodern}
\usepackage{textcomp}
\usepackage[a4paper,margin=2.5cm,headheight=14pt]{geometry}
\usepackage{amsmath,amssymb,amsthm}
\usepackage{amscd}  % KaTeX's CD diagrams; part of amsmath, so always there
\usepackage{mathtools}
\usepackage{bm}
\usepackage{graphicx}
\usepackage{xcolor}
\usepackage{array,booktabs,tabularx}
\usepackage{fancyvrb}
\usepackage[normalem]{ulem}
\usepackage{microtype}
% Optional, and only with their fonts: a small TeX installation has the
% packages without the fonts, which fails at the first formula.
\IfFileExists{mathrsfs.sty}{\IfFileExists{rsfs10.tfm}{\usepackage{mathrsfs}}{}}{}
\IfFileExists{stmaryrd.sty}{\IfFileExists{stmary10.tfm}{\usepackage{stmaryrd}}{}}{}
\IfFileExists{cancel.sty}{\usepackage{cancel}}{}
\usepackage[hyphens]{url}
\usepackage[pdfusetitle,bookmarksnumbered=false]{hyperref}
\usepackage{bookmark}
\definecolor{bplink}{HTML}{1F4E9C}
\definecolor{bpmuted}{HTML}{5F676E}
\hypersetup{colorlinks=true,linkcolor=bplink,urlcolor=bplink,citecolor=bplink}
\setlength{\parindent}{0pt}
\setlength{\parskip}{0.55em plus 0.2em minus 0.1em}
% Lists as tight as the page's.
\makeatletter
\def\@listi{\leftmargin\leftmargini\parsep 0pt\topsep 0.3em plus 0.1em\itemsep 0.1em plus 0.1em}
\let\@listI\@listi
\def\@listii{\leftmargin\leftmarginii\labelwidth\leftmarginii\advance\labelwidth-\labelsep\parsep 0pt\topsep 0.1em\itemsep 0.1em}
% An ordered list that does not start at 1.
\newcommand{\bpstartat}[1]{\setcounter{\@enumctr}{\numexpr#1-1\relax}}
\makeatother
\setcounter{tocdepth}{3}
\emergencystretch=3em
\allowdisplaybreaks

% A numbered statement, set as a paper sets one: "Theorem 2.1 (Title)."
% #1 is the lead ("Theorem~2.1"), #2 what follows it in roman.
\newenvironment{bpstatement}[2]{\par\addvspace{0.6\baselineskip}\noindent\textbf{#1}#2\textbf{.}\enspace\ignorespaces}
  {\par\addvspace{0.4\baselineskip}}
\newenvironment{bpstatementit}[2]{\begin{bpstatement}{#1}{#2}\itshape\ignorespaces}{\end{bpstatement}}
% A step: a relation with prose of its own, read under the object it starts from.
\newenvironment{bpstep}[1]{\par\addvspace{0.3\baselineskip}\begin{list}{}{\leftmargin=1.6em\rightmargin=0pt\topsep=0pt}\item\relax
  {\small\color{bpmuted}#1}\par\nopagebreak}{\end{list}\par}
% Notes about where an object is written out.
\newenvironment{bpnote}{\par\small\color{bpmuted}\itshape}{\par}
\newcommand{\bpomitted}[1]{\par{\small\color{bpmuted}\itshape #1\par}}
\newcommand{\bpbodyheading}[1]{\par\addvspace{0.4\baselineskip}\noindent\textbf{#1}\par\nopagebreak}
\newcommand{\bprule}{\par\medskip\noindent\hfil\rule{0.3\linewidth}{0.4pt}\hfil\par\medskip}
\newcommand{\bpboundary}[1]{\par\nopagebreak{\small\color{bpmuted}#1\par}\nopagebreak}
% A reference to something this export leaves out, and one that resolves nowhere.
\newcommand{\bpout}[1]{#1\textsuperscript{\hyperlink{bp-outside}{\ensuremath{\ast}}}}
\newcommand{\bpbroken}[1]{\textcolor{red!70!black}{[#1?]}}
\newcommand{\bpslug}[1]{\texttt{#1}}
% A bibliography entry: label, anchor, entry; the label hangs in the margin.
\newcommand{\bpbibitem}[3]{\par\setbox0\hbox{[#1]\enspace}\ifdim\wd0<5em\setbox0\hbox to 5em{[#1]\hfil}\fi
  \hangindent=\wd0 \hangafter=1 \noindent\box0\phantomsection\label{#2}#3\par\smallskip}
% A character the export knows no way to set.
\newcommand{\bpglyph}[1]{\fbox{\scriptsize #1}}
% Images, at their own size unless they are wider than the line.
\newsavebox{\bpimagebox}
\newcommand{\bpimage}[1]{\sbox{\bpimagebox}{\includegraphics{#1}}%
  \ifdim\wd\bpimagebox>\linewidth\resizebox{\linewidth}{!}{\usebox{\bpimagebox}}\else\usebox{\bpimagebox}\fi}
\newcommand{\bpmissingimage}[1]{\fbox{\parbox{0.8\linewidth}{\centering\small\color{bpmuted}Image not included: #1}}}
\DefineVerbatimEnvironment{bpverbatim}{Verbatim}{fontsize=\small,xleftmargin=1em}
\renewcommand{\qedsymbol}{\ensuremath{\blacksquare}}
\hypersetup{pdftitle={Weil Conjectures — 15 Pure sheaves on curves},pdfsubject={Chapter 15 Pure sheaves on curves; all levels},pdfcreator={blueprint}}
% KaTeX commands the website knows
\providecommand{\R}{\mathbb{R}}\providecommand{\N}{\mathbb{N}}\providecommand{\Z}{\mathbb{Z}}
\providecommand{\Q}{\mathbb{Q}}\providecommand{\C}{\mathbb{C}}
\providecommand{\Reals}{\mathbb{R}}\providecommand{\reals}{\mathbb{R}}\providecommand{\Complex}{\mathbb{C}}
\providecommand{\cnums}{\mathbb{C}}\providecommand{\natnums}{\mathbb{N}}\providecommand{\Bbb}[1]{\mathbb{#1}}
\providecommand{\bold}[1]{\mathbf{#1}}\providecommand{\frak}[1]{\mathfrak{#1}}\providecommand{\mathbbm}[1]{\mathbb{#1}}
\providecommand{\lang}{\langle}\providecommand{\rang}{\rangle}\providecommand{\bra}[1]{\langle #1\rvert}
\providecommand{\ket}[1]{\lvert #1\rangle}\providecommand{\braket}[1]{\langle #1\rangle}
\providecommand{\set}[1]{\{#1\}}\providecommand{\Set}[1]{\left\{#1\right\}}
\providecommand{\infin}{\infty}\providecommand{\isin}{\in}\providecommand{\plusmn}{\pm}\providecommand{\sub}{\subset}
\providecommand{\sube}{\subseteq}\providecommand{\supe}{\supseteq}\providecommand{\exist}{\exists}
\providecommand{\bull}{\bullet}\providecommand{\larr}{\leftarrow}\providecommand{\rarr}{\rightarrow}
\providecommand{\lrarr}{\leftrightarrow}\providecommand{\Larr}{\Leftarrow}\providecommand{\Rarr}{\Rightarrow}
\providecommand{\Lrarr}{\Leftrightarrow}\providecommand{\uarr}{\uparrow}\providecommand{\darr}{\downarrow}
\providecommand{\harr}{\leftrightarrow}\providecommand{\lArr}{\Leftarrow}\providecommand{\rArr}{\Rightarrow}
\providecommand{\hArr}{\Leftrightarrow}\providecommand{\alef}{\aleph}\providecommand{\alefsym}{\aleph}
\providecommand{\real}{\Re}\providecommand{\image}{\Im}\providecommand{\weierp}{\wp}
\providecommand{\Alpha}{\mathrm{A}}\providecommand{\Beta}{\mathrm{B}}\providecommand{\Epsilon}{\mathrm{E}}
\providecommand{\Zeta}{\mathrm{Z}}\providecommand{\Eta}{\mathrm{H}}\providecommand{\Iota}{\mathrm{I}}
\providecommand{\Kappa}{\mathrm{K}}\providecommand{\Mu}{\mathrm{M}}\providecommand{\Nu}{\mathrm{N}}
\providecommand{\Omicron}{\mathrm{O}}\providecommand{\Rho}{\mathrm{P}}\providecommand{\Tau}{\mathrm{T}}
\providecommand{\Chi}{\mathrm{X}}\providecommand{\omicron}{o}\providecommand{\gt}{>}\providecommand{\lt}{<}
\providecommand{\And}{\&}\providecommand{\hdashline}{\hline}\providecommand{\operatornamewithlimits}[1]{\operatorname*{#1}}
\providecommand{\llbracket}{[\![}\providecommand{\rrbracket}{]\!]}\providecommand{\cancel}[1]{#1}
\providecommand{\mathscr}{\mathcal}
% the characters beyond ASCII this document uses
\DeclareUnicodeCharacter{2161}{\textup{II}}
\DeclareUnicodeCharacter{2162}{\textup{III}}
% the project's macros ([katex.macros] in blueprint.toml)
\providecommand{\A}{}\renewcommand{\A}{\mathbf A}
\providecommand{\AEn}{}\renewcommand{\AEn}{A_{E,n}}
\providecommand{\Aaf}{}\renewcommand{\Aaf}{\mathbf A}
\providecommand{\Ab}{}\renewcommand{\Ab}{\mathsf{Ab}}
\providecommand{\AffSch}{}\renewcommand{\AffSch}{\operatorname{AffSch}}
\providecommand{\Aone}{}\renewcommand{\Aone}{\mathbf A^1}
\providecommand{\Bl}{}\renewcommand{\Bl}{\operatorname{Bl}}
\providecommand{\Br}{}\renewcommand{\Br}{\mathrm{Br}}
\providecommand{\C}{}\renewcommand{\C}{\mathbb C}
\providecommand{\Cech}{}\renewcommand{\Cech}{\check{H}}
\providecommand{\Coext}{}\renewcommand{\Coext}{\operatorname{Coext}}
\providecommand{\Cone}{}\renewcommand{\Cone}{\operatorname{Cone}}
\providecommand{\Cons}{}\renewcommand{\Cons}{\operatorname{Cons}}
\providecommand{\Cref}{}\renewcommand{\Cref}[1]{\text{#1}}
\providecommand{\D}{}\renewcommand{\D}{D}
\providecommand{\DLAFT}{}\renewcommand{\DLAFT}{\operatorname{FT}}
\providecommand{\DLAFr}{}\renewcommand{\DLAFr}{\operatorname{Fr}}
\providecommand{\DLAQ}{}\renewcommand{\DLAQ}{{\overline{\mathbf Q}_\ell}}
\providecommand{\DLRQ}{}\renewcommand{\DLRQ}{{\overline{\mathbf Q}_\ell}}
\providecommand{\Dbc}{}\renewcommand{\Dbc}{D^b_c}
\providecommand{\Dc}{}\renewcommand{\Dc}{D_{\mathrm{cons}}}
\providecommand{\Dctf}{}\renewcommand{\Dctf}{D^b_{ctf}}
\providecommand{\Dic}{}\renewcommand{\Dic}{D_{\IndCons}}
\providecommand{\Dplus}{}\renewcommand{\Dplus}{D^{+}}
\providecommand{\End}{}\renewcommand{\End}{\operatorname{End}}
\providecommand{\Et}{}\renewcommand{\Et}{\mathrm{\acute{e}t}}
\providecommand{\Ex}{}\renewcommand{\Ex}{\mathrm{Ex}}
\providecommand{\Ext}{}\renewcommand{\Ext}{\operatorname{Ext}}
\providecommand{\F}{}\renewcommand{\F}{\mathbb F}
\providecommand{\FT}{}\renewcommand{\FT}{\operatorname{FT}}
\providecommand{\Fbar}{}\renewcommand{\Fbar}{{\overline{\mathbf F}_q}}
\providecommand{\Fell}{}\renewcommand{\Fell}{\mathbf F_\ell}
\providecommand{\Fib}{}\renewcommand{\Fib}{\operatorname{Fib}}
\providecommand{\Fix}{}\renewcommand{\Fix}{\operatorname{Fix}}
\providecommand{\Fp}{}\renewcommand{\Fp}{{\mathbf F_{p}}}
\providecommand{\Fq}{}\renewcommand{\Fq}{{\mathbf F_q}}
\providecommand{\Fqn}{}\renewcommand{\Fqn}{{\mathbf F_{q^{n}}}}
\providecommand{\Frob}{}\renewcommand{\Frob}{\operatorname{Frob}}
\providecommand{\GL}{}\renewcommand{\GL}{\operatorname{GL}}
\providecommand{\Gal}{}\renewcommand{\Gal}{\operatorname{Gal}}
\providecommand{\Hens}{}\renewcommand{\Hens}{\operatorname{Hens}}
\providecommand{\Hilb}{}\renewcommand{\Hilb}{\operatorname{Hilb}}
\providecommand{\Hom}{}\renewcommand{\Hom}{\operatorname{Hom}}
\providecommand{\IndCons}{}\renewcommand{\IndCons}{\operatorname{IndCons}}
\providecommand{\Kar}{}\renewcommand{\Kar}{\operatorname{Kar}}
\providecommand{\LT}{}\renewcommand{\LT}{\operatorname{LT}}
\providecommand{\Lie}{}\renewcommand{\Lie}{\operatorname{Lie}}
\providecommand{\Loc}{}\renewcommand{\Loc}{\operatorname{Loc}}
\providecommand{\Mod}{}\renewcommand{\Mod}{\mathsf{Mod}}
\providecommand{\N}{}\renewcommand{\N}{\mathbb{N}}
\providecommand{\Needspace}{}\renewcommand{\Needspace}{\par}
\providecommand{\O}{}\renewcommand{\O}{{\mathcal{O}}}
\providecommand{\Ob}{}\renewcommand{\Ob}{\mathrm{Ob}}
\providecommand{\Oe}{}\renewcommand{\Oe}{\mathcal O_E}
\providecommand{\PP}{}\renewcommand{\PP}{\mathbf{P}}
\providecommand{\Perf}{}\renewcommand{\Perf}{\operatorname{Perf}}
\providecommand{\Pic}{}\renewcommand{\Pic}{\mathrm{Pic}}
\providecommand{\Pm}{}\renewcommand{\Pm}{\mathbf P^{m}}
\providecommand{\Pone}{}\renewcommand{\Pone}{\mathbf P^{1}}
\providecommand{\Pqn}{}\renewcommand{\Pqn}{\mathbf P^{n}}
\providecommand{\ProFin}{}\renewcommand{\ProFin}{\operatorname{ProFin}}
\providecommand{\Proet}{}\renewcommand{\Proet}{\mathrm{pro\acute{e}t}}
\providecommand{\Q}{}\renewcommand{\Q}{\mathbb Q}
\providecommand{\Qb}{}\renewcommand{\Qb}{\overline{\mathbf Q}_{\ell}}
\providecommand{\Qbar}{}\renewcommand{\Qbar}{\overline{\mathbf Q}}
\providecommand{\QbarX}{}\renewcommand{\QbarX}{\overline{\mathbf Q}_{\ell,X}}
\providecommand{\Qbarl}{}\renewcommand{\Qbarl}{{\overline{\mathbf Q}_{\ell}}}
\providecommand{\Qell}{}\renewcommand{\Qell}{{\mathbf{Q}_{\ell}}}
\providecommand{\Ql}{}\renewcommand{\Ql}{\mathbf Q_{\ell}}
\providecommand{\Qlb}{}\renewcommand{\Qlb}{{\overline{\mathbf Q}_\ell}}
\providecommand{\Qlbar}{}\renewcommand{\Qlbar}{\overline{\mathbf Q}_\ell}
\providecommand{\R}{}\renewcommand{\R}{\mathbb{R}}
\providecommand{\RG}{}\renewcommand{\RG}{\operatorname{\mathrm{R}\Gamma}}
\providecommand{\RGamma}{}\renewcommand{\RGamma}{R\Gamma}
\providecommand{\RGammac}{}\renewcommand{\RGammac}{R\Gamma_c}
\providecommand{\RHom}{}\renewcommand{\RHom}{R\mathcal{H}om}
\providecommand{\RHomX}{}\renewcommand{\RHomX}{R\mathcal H om_X}
\providecommand{\Rat}{}\renewcommand{\Rat}{\operatorname{IsRat}}
\providecommand{\Res}{}\renewcommand{\Res}{\operatorname{Res}}
\providecommand{\Rlim}{}\renewcommand{\Rlim}{R\!\varprojlim}
\providecommand{\SIm}{}\renewcommand{\SIm}{\operatorname{SIm}}
\providecommand{\Sch}{}\renewcommand{\Sch}{\mathsf{Sch}}
\providecommand{\Sh}{}\renewcommand{\Sh}{\operatorname{Sh}}
\providecommand{\Span}{}\renewcommand{\Span}{\operatorname{span}}
\providecommand{\Spec}{}\renewcommand{\Spec}{\operatorname{Spec}}
\providecommand{\Tr}{}\renewcommand{\Tr}{\operatorname{Tr}}
\providecommand{\Trsm}{}\renewcommand{\Trsm}{\operatorname{Tr}^{\mathrm{sm}}}
\providecommand{\Wt}{}\renewcommand{\Wt}{\operatorname{Wt}}
\providecommand{\Z}{}\renewcommand{\Z}{\mathbb{Z}}
\providecommand{\Zar}{}\renewcommand{\Zar}{\mathrm{zar}}
\providecommand{\Zell}{}\renewcommand{\Zell}{{\mathbf{Z}_{\ell}}}
\providecommand{\Zl}{}\renewcommand{\Zl}{{\mathbf Z_\ell}}
\providecommand{\affet}{}\renewcommand{\affet}[1]{{#1}^{\mathrm{aff}}_{\mathrm{ét}}}
\providecommand{\affproet}{}\renewcommand{\affproet}[1]{{#1}^{\mathrm{aff}}_{\mathrm{proét}}}
\providecommand{\boxt}{}\renewcommand{\boxt}{\boxtimes_A^{\mathrm L}}
\providecommand{\cF}{}\renewcommand{\cF}{\mathcal{F}}
\providecommand{\cHom}{}\renewcommand{\cHom}{R\mathcal{H}\!om}
\providecommand{\cI}{}\renewcommand{\cI}{\mathcal{I}}
\providecommand{\cL}{}\renewcommand{\cL}{\mathcal L}
\providecommand{\cM}{}\renewcommand{\cM}{\mathcal M}
\providecommand{\cN}{}\renewcommand{\cN}{\mathcal N}
\providecommand{\cO}{}\renewcommand{\cO}{\mathcal O}
\providecommand{\cT}{}\renewcommand{\cT}{\mathcal T}
\providecommand{\checked}{}\renewcommand{\checked}{\color{red}\aftergroup\restorecolor}
\providecommand{\coeff}{}\renewcommand{\coeff}{\operatorname{coeff}}
\providecommand{\coker}{}\renewcommand{\coker}{\operatorname{coker}}
\providecommand{\colim}{}\renewcommand{\colim}{\operatorname{colim}}
\providecommand{\constructedby}{}\renewcommand{\constructedby}[1]{\par\smallskip\noindent\textcolor{purple!60!black}{\textsf{Constructed by:} }
  \footnotesize\sloppy\nolinkurl{#1}\par}
\providecommand{\deepresult}{}\renewcommand{\deepresult}[1]{\par\smallskip\noindent\textcolor{orange!45!black}{
  \textsf{Blueprint role: named deep theorem node; see the cluster proof architecture.}}
  \par\noindent\textcolor{orange!45!black}{\textsf{Mathematical source:} }#1\par}
\providecommand{\dt}{}\renewcommand{\dt}{\mathbin{\otimes^{\mathbf L}}}
\providecommand{\dw}{}\renewcommand{\dw}{\operatorname{dw}}
\providecommand{\et}{}\renewcommand{\et}[1]{{#1}_{\mathrm{ét}}}
\providecommand{\ev}{}\renewcommand{\ev}{\operatorname{ev}}
\providecommand{\fib}{}\renewcommand{\fib}{\operatorname{fib}}
\providecommand{\formalizationnote}{}\renewcommand{\formalizationnote}[1]{\begin{quote}\small\raggedright\textbf{Formalization note.} #1\end{quote}}
\providecommand{\fpqc}{}\renewcommand{\fpqc}{X_{\mathrm{fpqc}}}
\providecommand{\gammacont}{}\renewcommand{\gammacont}{\Gamma_{\mathrm{cont}}}
\providecommand{\gm}{}\renewcommand{\gm}{\mathbf G_m}
\providecommand{\hocolim}{}\renewcommand{\hocolim}{\operatorname{hocolim}}
\providecommand{\holim}{}\renewcommand{\holim}{\operatorname{holim}}
\providecommand{\hrB}{}\renewcommand{\hrB}{\@ifnextchar]{\hrEnd}{\hrStep}}
\providecommand{\hrEnd}{}\renewcommand{\hrEnd}{\@firstoftwo{}}
\providecommand{\hrStep}{}\renewcommand{\hrStep}{\@firstoftwo{\hrB}}
\providecommand{\id}{}\renewcommand{\id}{\operatorname{id}}
\providecommand{\kunFl}{}\renewcommand{\kunFl}{\mathbf F_\ell}
\providecommand{\kunZn}{}\renewcommand{\kunZn}{\mathbf Z/n\mathbf Z}
\providecommand{\kuncd}{}\renewcommand{\kuncd}{\operatorname{cd}}
\providecommand{\kundim}{}\renewcommand{\kundim}{\dim}
\providecommand{\kunotimes}{}\renewcommand{\kunotimes}[1]{\otimes^{\mathbf L}_{#1}}
\providecommand{\kunsupp}{}\renewcommand{\kunsupp}{\operatorname{Supp}}
\providecommand{\lotimes}{}\renewcommand{\lotimes}{\otimes^{\mathrm L}}
\providecommand{\m}{}\renewcommand{\m}{{\mathfrak{m}}}
\providecommand{\n}{}\renewcommand{\n}{{\mathfrak{n}}}
\providecommand{\nagataGap}{}\renewcommand{\nagataGap}[1]{\par\noindent\textbf{Formalization gap.} #1\par}
\providecommand{\openobligation}{}\renewcommand{\openobligation}[1]{\par\smallskip\noindent\textcolor{red!65!black}{\textsf{Open obligation:} }
  \footnotesize\sloppy\nolinkurl{#1}\par}
\providecommand{\p}{}\renewcommand{\p}{{\mathfrak{p}}}
\providecommand{\pr}{}\renewcommand{\pr}{\operatorname{pr}}
\providecommand{\preshv}{}\renewcommand{\preshv}[1]{\mathrm{PreShv}(#1)}
\providecommand{\pro}{}\renewcommand{\pro}[1]{\mathrm{Pro}(#1)}
\providecommand{\proet}{}\renewcommand{\proet}{X_{\mathrm{proét}}}
\providecommand{\q}{}\renewcommand{\q}{{\mathfrak{q}}}
\providecommand{\rad}{}\renewcommand{\rad}{\operatorname{rad}}
\providecommand{\rank}{}\renewcommand{\rank}{\operatorname{rank}}
\providecommand{\red}{}\renewcommand{\red}{\color{red}}
\providecommand{\restorecolor}{}\renewcommand{\restorecolor}{\color{black}}
\providecommand{\retiredobligation}{}\renewcommand{\retiredobligation}[1]{\par\smallskip\noindent\textcolor{red!75!black}{\textsf{Retired invalid route:} }
  \footnotesize\sloppy\nolinkurl{#1}\par}
\providecommand{\rk}{}\renewcommand{\rk}{\operatorname{rk}}
\providecommand{\sep}{}\renewcommand{\sep}{\textnormal{sep}}
\providecommand{\shv}{}\renewcommand{\shv}[1]{\mathrm{Shv}(#1)}
\providecommand{\spec}{}\renewcommand{\spec}[1]{\mathrm{Spec}(#1)}
\providecommand{\stacks}{}\renewcommand{\stacks}[1]{\href{https://stacks.math.columbia.edu/tag/#1}{\textnormal{Stacks, Tag~#1}}}
\providecommand{\stacksproject}{}\renewcommand{\stacksproject}[1]{{\cite[\href{https://stacks.math.columbia.edu/tag/#1}{Tag #1}]{stacks-project}}}
\providecommand{\statusnote}{}\renewcommand{\statusnote}[1]{\begin{quote}\small\raggedright\textbf{Status.} #1\end{quote}}
\providecommand{\tildeideal}{}\renewcommand{\tildeideal}[2]{\widetilde{(#1, #2)}}
\providecommand{\tildering}{}\renewcommand{\tildering}[3]{{#1}_{\widetilde{(#2, #3)}}}
\providecommand{\tr}{}\renewcommand{\tr}{\operatorname{tr}}
\providecommand{\wt}{}\renewcommand{\wt}{\operatorname{wt}}

\begin{document}
\pagestyle{plain}
\begin{center}
{\LARGE\bfseries Weil Conjectures\par}
\vspace{0.7em}
{\large Chapter 15: Pure sheaves on curves\par}
\vspace{0.5em}
{\small\color{bpmuted}Snapshot 2026-09-26, revision 8e4cb88e2a\par}
\end{center}
\vspace{0.5em}
{\small\tableofcontents}
\bigskip\hrule\bigskip

% object: chap-Pure-sheaves-on-curves (blueprint/chap-Pure-sheaves-on-curves/_section.md)
\section*{\texorpdfstring{15\quad Pure sheaves on curves}{15 Pure sheaves on curves}}
\phantomsection\label{bp1}%
\addcontentsline{toc}{section}{\texorpdfstring{15\quad Pure sheaves on curves}{15 Pure sheaves on curves}}
% text: prose-15-01
\phantomsection\label{bp2}%
How much does the trace function of a sheaf know about its Frobenius eigenvalues? For a sheaf of rank one, pure of weight \(w\), on a dense open \(U_0\) of the line, it knows the weight. Indeed, every trace over \(\mathbf F_{q^n}\) has absolute value \(q^{nw/2}\), so the sum \(S_n\) of the \(|t(x)|^2\) over \(U_0(\mathbf F_{q^n})\) is \(q^{nw}\#U_0(\mathbf F_{q^n})\), about \(q^{n(w+1)}\), and the growth rate \(q^{w+1}\) of these sums therefore determines \(w\). In higher rank, however, a trace is a sum of eigenvalues, and the eigenvalues can cancel. For example, on \(\mathcal B_b\oplus\mathcal B_{-b}\) the trace at a point of degree \(n\) is \(b^n+(-b)^n\), which is zero for every odd \(n\). We show that the cancellation cannot persist: for a nonzero pure lisse sheaf of weight \(w\) on an affine curve, \(\limsup_n\frac1n\log_qS_n=w+1\) (Theorem \hyperref[bp41]{15.14}).

The reason is positivity once more. Recall that a sheaf \(\mathcal G\) is real if the characteristic polynomials of Frobenius on its stalks have real coefficients under \(\iota\) (Definition \bpout{14.21}). For such \(\mathcal G\), pure of weight \(w\), \(|t_{\mathcal G}|^2\) is the trace function of \(\mathcal G\otimes\mathcal G\), and hence the generating series of the \(S_n\) is the logarithmic derivative of \(L(\mathcal G\otimes\mathcal G,t)\). Positivity of the Euler factors (Lemma \bpout{14.22}) shows that \(L(\mathcal G\otimes\mathcal G,t)\) is not a polynomial, and Lemma \bpout{14.20} puts its poles on \(|t|=q^{-w-1}\). Together these force the sums \(S_n\) to reach the size \(q^{n(w+1)}\) infinitely often. The restriction to real sheaves costs nothing, since an arbitrary pure sheaf becomes real after a twist to weight zero and the addition of its dual (Proposition \hyperref[bp27]{15.9}).

This is what the Fourier argument of Chapter \bpout{17} needs. That argument starts from a pure lisse sheaf \(M\) of weight \(w\) on a dense open \(j:U_0\hookrightarrow\mathbf A^1\) and produces a pure lisse sheaf \(M'\) on a dense open \(j'\) of the dual line. It knows the trace function of \(M'\) exactly, as the discrete Fourier transform of that of \(M\), but it cannot see the weight of \(M'\) directly. Since Plancherel's formula multiplies the sums \(S_n\) by \(q^n\), the theorem shows that \(M'\) has weight \(w+1\).

The group we finally want to bound, \(H^1_c(\mathbf A^1,j_*M)\), is the stalk of \(j'_*M'\) at the origin. But the origin may be one of the points where \(M'\) is not lisse, and there we need to know that \(j'_*\) does not create larger weights (Proposition \hyperref[bp13]{15.5}).

% object: sec-Weights-of-middle-extensions (blueprint/chap-Pure-sheaves-on-curves/sec-Weights-of-middle-extensions/_section.md), content hidden
\subsection*{\texorpdfstring{15.1\quad Weights of middle extensions}{15.1 Weights of middle extensions}}
\phantomsection\label{bp3}%
\addcontentsline{toc}{subsection}{\texorpdfstring{15.1\quad Weights of middle extensions}{15.1 Weights of middle extensions}}
% text: prose-15-02
\phantomsection\label{bp4}%
Let \(j\) be the inclusion of a dense open in a smooth curve, and let \(\mathcal F\) be a lisse sheaf on the open. The Fourier argument needs an extension of \(\mathcal F\) to the whole curve that loses nothing at the boundary points. The ordinary direct image \(j_*\mathcal F\) is such an extension, whereas extension by zero loses something. For example, for the constant sheaf on \(\mathbf G_m\subset\mathbf A^1\), \(j_!\overline{\mathbf Q}_\ell\) has stalk \(0\) at the origin, while the constant sheaf on \(\mathbf A^1\) is the natural extension. In general, by Lemma \hyperref[bp6]{15.2}, \(j_*\mathcal F\) recovers a lisse extension whenever one exists, and its stalk at a boundary point is the space of inertia invariants.

If \(\mathcal F\) is punctually of weights \(\le a\), then so is \(j_*\mathcal F\), at the boundary points as well (Proposition \hyperref[bp13]{15.5}). In Chapter \bpout{17} this bounds the weights of the boundary terms in the Fourier transform of a middle extension.

% object: definition-middle-extension (blueprint/chap-Pure-sheaves-on-curves/sec-Weights-of-middle-extensions/definition-middle-extension/_section.md)
\phantomsection\label{bp5}%
\begin{bpstatement}{Definition~15.1}{ (Middle extension)}
Let \(j:U\hookrightarrow X\) be a dense open in a smooth curve and let \(\mathcal F\) be a lisse sheaf on \(U\). The \emph{middle extension} of \(\mathcal F\) is the ordinary direct image \(j_*\mathcal F\) of sheaves.
\end{bpstatement}

% object: lemma-boundary-stalk (blueprint/chap-Pure-sheaves-on-curves/sec-Weights-of-middle-extensions/lemma-boundary-stalk/_section.md)
\phantomsection\label{bp6}%
\begin{bpstatementit}{Lemma~15.2}{ (The boundary stalk)}
Let \(j:U_0\hookrightarrow X_0\) be a dense open in a smooth curve, let \(\mathcal F\) be a lisse Weil sheaf on \(U_0\), and let \(j_*\) be the ordinary direct image of sheaves. Let \(\bar z\) be a geometric point over a point \(z\) of \(X_0-U_0\), and let \(I_z\) be the inertia group at \(z\), acting on the local generic fibre \(V\) of \(\mathcal F\). There is a canonical Frobenius-equivariant identification

\[
(j_*\mathcal F)_{\bar z}=V^{I_z}.
\]

If \(\mathcal F\) extends to a lisse sheaf on \(X_0\), the adjunction map from that extension to \(j_*\mathcal F\) is an isomorphism.
\end{bpstatementit}

% proof: proof/lemma-boundary-stalk
\begin{proof}[{Proof}]\phantomsection\label{bp7}%
We prove the formula for finite coefficients, where it is Galois theory of the punctured henselian neighbourhood, and pass to the limit.

The stalk of \(j_*\mathcal F\) at \(\bar z\) is the colimit of the sections of \(\mathcal F\) over the inverse images of \(U\) in étale neighbourhoods of \(\bar z\), that is, the group of sections over the punctured spectrum of the strict henselization at \(\bar z\). This punctured spectrum is the spectrum of a field, whose absolute Galois group is the inertia group \(I_z\), so finite locally constant sheaves there correspond to finite continuous representations of the inertia group, and their sections are the invariant vectors. Applying this to the finite reductions of a stable lattice and passing to the inverse limit, we obtain the formula with integral coefficients.

Two compatibilities take us to rational coefficients and to \(\overline{\mathbf Q}_\ell\). Rationalization commutes with taking invariants, since a power of the uniformizer carries an invariant rational vector into the lattice.  A finite extension of scalars also commutes with invariants. Indeed, it is exact, hence commutes with finite intersections of kernels, and in a finite-dimensional space the intersection of the kernels of all \(g-1\) is the intersection of finitely many of them. This gives the formula with rational coefficients, after finite extension of scalars, and hence also over \(\overline{\mathbf Q}_\ell\). Since all these maps commute with the action of the decomposition group, the identification is Frobenius equivariant. Finally, if \(\mathcal F\) extends to a lisse sheaf on \(X_0\), the local inertia acts trivially, and so the adjunction map is an isomorphism on every geometric stalk.
\end{proof}

% object: example-boundary-stalks-g-m (blueprint/chap-Pure-sheaves-on-curves/sec-Weights-of-middle-extensions/example-boundary-stalks-g-m/_section.md)
\phantomsection\label{bp8}%
\begin{bpstatement}{Example~15.3}{ (Boundary stalks on $\mathbf G_m\subset\mathbf A^1$)}
Let \(j:\mathbf G_{m,k_0}\hookrightarrow\mathbf A^1_{k_0}\) and \(z=0\). The covers \(u^n=t\) are totally ramified at \(0\), so the inertia \(I_0\) maps onto the monodromy of a Kummer sheaf. For the Kummer sheaf of a nontrivial character, \(I_0\) acts on the line \(V\) through that character. Hence \(V^{I_0}=0\), and the stalk of \(j_*\) of this sheaf at \(\bar0\) is zero, so here \(j_*\) and \(j_!\) agree. For the Kummer extension \(\mathcal E\) of Example \bpout{13.4}, on the other hand, \(I_0\) acts by the matrices \(\bigl(\begin{smallmatrix}1&\kappa\\0&1\end{smallmatrix}\bigr)\) with \(\kappa\) running over \(\mathbf Z_\ell(1)\), whose common fixed vectors form the line \(\overline{\mathbf Q}_\ell(1)\). So the stalk of \(j_*\mathcal E\) at \(0\) is \(\overline{\mathbf Q}_\ell(1)\), of weight \(-2\), while \(\mathcal E\) has weights \(-2\) and \(0\); the stalk at the boundary keeps only part of the weights.
\end{bpstatement}

% text: prose-15-03
\phantomsection\label{bp9}%
To bound \(j_*\mathcal F\) we compare the \(L\)-functions of \(\mathcal F\) and of \(j_*\mathcal F\) on a disc. To see that the Euler product converges there, we need a crude bound for the number of points of a curve.

% object: lemma-point-count-bound (blueprint/chap-Pure-sheaves-on-curves/sec-Weights-of-middle-extensions/lemma-point-count-bound/_section.md)
\phantomsection\label{bp10}%
\begin{bpstatementit}{Lemma~15.4}{ (A point-count bound)}
For a smooth affine curve \(U_0/k_0\) there is a constant \(D\) such that \(\#U_0(k_n)\le Dq^n\) for all \(n\ge1\).  For a finite \(k_0\)-scheme, the number of \(k_n\)-points is bounded independently of \(n\).
\end{bpstatementit}

% proof: proof/lemma-point-count-bound
\begin{proof}[{Proof}]\phantomsection\label{bp11}%
We map \(U_0\) finitely and flatly onto the affine line, whose number of \(k_n\)-points is \(q^n\), and bound the size of the fibres.

By Noether normalization [Stacks, \href{https://stacks.math.columbia.edu/tag/00OY}{00OY}] there is a finite map \(U_0\to\mathbf A^1_{k_0}\), and the coordinate ring of \(U_0\) is a finite torsion-free module over \(k_0[t]\). To see this, let \(f\ne0\) in \(k_0[t]\) and \(fa=0\). On each irreducible component of \(U_0\) the image of \(f\) is nonzero, since otherwise the component would be finite over \(k_0[t]/(f)\) and hence of dimension zero. Since the coordinate ring of each component is a domain, \(fa=0\) forces \(a\) to vanish on every component, and then \(a=0\) because \(U_0\) is reduced.  As \(k_0[t]\) is a principal ideal domain, it follows that the coordinate ring of \(U_0\) is free of some rank \(D\).  Hence, after any extension of the base field, every fibre has a coordinate ring of dimension \(D\), and by the Chinese remainder theorem it has at most \(D\) rational points.  Summing over \(\mathbf A^1(k_n)\), we obtain the first bound.  The second holds because the number of rational points of a finite scheme is at most its number of geometric points.
\end{proof}

% text: prose-15-04
\phantomsection\label{bp12}%
In the next proof we use the lemma to show that the logarithm of the Euler product of \(\mathcal F\) converges on a disc of radius \(q^{-a/2-1}\), and on this disc we compare it with \(j_*\mathcal F\).

% object: proposition-weights-j-boundary-points (blueprint/chap-Pure-sheaves-on-curves/sec-Weights-of-middle-extensions/proposition-weights-j-boundary-points/_section.md)
\phantomsection\label{bp13}%
\begin{bpstatementit}{Proposition~15.5}{ (Weights of $j_*$ at the boundary points)}
Let \(j:U_0\hookrightarrow X_0\) be a dense open of a smooth geometrically connected curve.  If a lisse Weil sheaf \(\mathcal F\) on \(U_0\) is punctually of weights \(\le a\), so is \(j_*\mathcal F\) on \(X_0\).
\end{bpstatementit}

% proof: proof/proposition-weights-j-boundary-points
\begin{proof}[{Proof}]\phantomsection\label{bp14}%
Comparing the \(L\)-functions of \(\mathcal F\) and of \(j_*\mathcal F\) gives the weaker bound \(a+2\) at the boundary, and tensor powers remove the extra \(2\).

The assertion is local at a point of \(X_0-U_0\), so we may replace \(X_0\) by an affine neighbourhood of that point. The complement of \(U_0\) is an effective Cartier divisor, so after shrinking this neighbourhood it is cut out by one equation \(f\), and then \(U_0=D(f)\) is affine as well. Write \(i:Z_0=X_0-U_0\hookrightarrow X_0\). The hypotheses are preserved by this replacement. Indeed, forming \(j_*\) commutes with restriction to open subschemes, and the restriction of \(\mathcal F\) is still punctually of weights \(\le a\). Moreover, a nonempty open of \(X_0\) is still geometrically connected.

The boundary Euler factors of \(j_*\mathcal F\) are the ratio \(L_X/L_U\) of the two \(L\)-functions defined below. We show that \(L_U\) has no zero and \(L_X\) no pole on the disc \(|t|<q^{-a/2-1}\), so that the boundary factors have no pole there, and this bounds their eigenvalues.

\emph{The \(L\)-function on \(U_0\).} Put \(L_U(t)=L(U_0,\mathcal F,t)\). We show that \(\iota L_U\) is holomorphic and nonvanishing on the disc, by bounding the coefficients of its logarithm. The logarithm of the Euler product has coefficients

\[
\log\iota L_U(t)=\sum_{n\ge1}\frac{c_n}{n}t^n,
\qquad
c_n=\sum_{d\mid n}d\sum_{\deg(x)=d}
\iota\operatorname{Tr}(\operatorname{Frob}_x^{n/d}\mid\mathcal F_{\bar x}).
\]

Let \(r=\operatorname{rank}\mathcal F\). For a closed point \(x\) of degree \(d\mid n\), the bound on the weights gives \(|\iota\operatorname{Tr}(\operatorname{Frob}_x^{n/d}\mid\mathcal F_{\bar x})|
\le rq_x^{(n/d)a/2}=rq^{na/2}\). Since \(\sum_{d\mid n}d\cdot\#\{x\in|U_0|:\deg(x)=d\}=\#U_0(k_n)\), Lemma \hyperref[bp10]{15.4} gives \(|c_n|\le rDq^{n(a/2+1)}\). Hence the logarithm converges on the disc \(|t|<R=q^{-a/2-1}\), and its exponential is a nonvanishing holomorphic function there. By the trace formula, its Taylor series is that of the rational function \(\iota L_U(t)\), so the two agree on the whole disc.

\emph{The \(L\)-function on \(X_0\).} Let \(L_X(t)=L(X_0,j_*\mathcal F,t)\). We show that \(\iota L_X\) is holomorphic on the same disc. By the trace formula its poles come from \(H^0_c\) and \(H^2_c\) of \(j_*\mathcal F\), and we compute both groups.

The group \(H^0_c\) vanishes. Indeed, the support of a compactly supported section of \(j_*\mathcal F\) on \(X\) is closed and proper over \(k\), and since \(X\) is an affine curve, the support is finite. So the section vanishes on a dense open subset of \(U\), and as \(\mathcal F\) is lisse and \(U\) is connected, it vanishes on all of \(U\). Since a section of \(j_*\mathcal F\) is determined by its restriction to \(U\), the section is zero. For \(H^2_c\), note that the quotient sheaf in

\[
0\longrightarrow j_!\mathcal F\longrightarrow j_*\mathcal F
\longrightarrow i_*i^*j_*\mathcal F\longrightarrow0
\]

is supported on the finite boundary, whose cohomology is concentrated in degree zero.  From the long exact sequence we obtain \(H_c^2(X,j_*\mathcal F)=H_c^2(U,\mathcal F)\), and by Lemma \bpout{14.20} its eigenvalues have absolute value at most \(q^{a/2+1}\).  By the trace formula, every pole of \(\iota L_X\) is the inverse of a Frobenius eigenvalue on \(H^0_c\) or \(H^2_c\). Since \(H^0_c=0\), the poles have absolute value at least \(q^{-a/2-1}=R\). Hence \(\iota L_X\) is holomorphic on \(|t|<R\).

\emph{The boundary factors.} On \(U_0\) the sheaves \(\mathcal F\) and \(j_*\mathcal F\) have the same stalks, so the two Euler products differ only by the factors at the points of \(Z_0\), and we get

\[
L_X(t)=L_U(t)\prod_{z\in Z_0}
\det(1-\operatorname{Frob}_z t^{\deg(z)}\mid(j_*\mathcal F)_{\bar z})^{-1}.
\]

Since \(L_U\) has no zero on the disc, the finite product on the right is holomorphic there. As it is the reciprocal of a product of polynomials with constant term one, none of these polynomials has a zero on the disc.  Now, for an eigenvalue \(\alpha\) at \(z\), the roots of \(1-\iota(\alpha)t^{\deg(z)}\) have absolute value \(|\iota(\alpha)|^{-1/\deg(z)}\). This number is therefore at least \(R\), which is the preliminary bound

\[
|\iota(\alpha)|^2\le q_z^{a+2}.
\]

\emph{Tensor powers.} The error \(2\) in the exponent comes from the twist \((-1)\) on \(H^2_c\). Since it does not grow with the weight, it becomes negligible for high tensor powers. We therefore apply the preliminary bound to \(\mathcal F^{\otimes m}\) for every \(m\ge1\), a sheaf punctually of weights \(\le ma\).  Let \(V\) be the local generic fibre of \(\mathcal F\) at \(z\). By Lemma \hyperref[bp6]{15.2}, the eigenvalue \(\alpha\) of \(\operatorname{Frob}_z\) on \((j_*\mathcal F)_{\bar z}\) has an eigenvector \(v\in V^{I_z}\), and then \(v^{\otimes m}\) is a nonzero invariant eigenvector in \(V^{\otimes m}\) with eigenvalue \(\alpha^m\). By the same lemma, \(\alpha^m\) is therefore an eigenvalue on the stalk of \(j_*(\mathcal F^{\otimes m})\) at \(\bar z\), and the preliminary bound for \(\mathcal F^{\otimes m}\) gives \(|\iota(\alpha)|^{2m}\le q_z^{ma+2}\).  Taking \(m\)th roots and letting \(m\) tend to infinity, we obtain \(|\iota(\alpha)|^2\le q_z^a\) at every boundary point. On \(U_0\), finally, the bound is the hypothesis.
\end{proof}

% text: prose-15-05
\phantomsection\label{bp15}%
The proposition gives only an upper bound at the boundary, and in general nothing better is true. For example, for a family of elliptic curves with multiplicative reduction at \(z\), the first cohomology of the fibres is pure of weight \(1\), while its inertia invariants at \(z\) form a line of weight \(0\). When the inertia acts trivially, however, we get the lower bound by duality.

% object: corollary-unramified-boundary-point (blueprint/chap-Pure-sheaves-on-curves/sec-Weights-of-middle-extensions/corollary-unramified-boundary-point/_section.md)
\phantomsection\label{bp16}%
\begin{bpstatementit}{Corollary~15.6}{ (An unramified boundary point)}
Let \(\mathcal F\) be a lisse sheaf, pure of weight \(a\), on a smooth geometrically connected curve with a point removed. If the inertia at the removed point acts trivially, then the stalk at that point of the direct image of \(\mathcal F\) to the curve is pure of weight \(a\). In particular, a lisse Weil sheaf on a smooth geometrically connected curve that is pure on a dense open is pure on the whole curve.
\end{bpstatementit}

% proof: proof/corollary-unramified-boundary-point
\begin{proof}[{Proof}]\phantomsection\label{bp17}%
Purity follows from the upper bounds for the sheaf and for its dual, provided the stalk of the dual is the dual of the stalk, and trivial inertia guarantees this. So we apply Proposition \hyperref[bp13]{15.5} to the sheaf and to its dual, which is pure of weight \(-a\).  Since the inertia acts trivially, the stalk of the extended dual is the dual of the stalk of the extension, and so every eigenvalue \(\alpha\) satisfies \(|\iota\alpha|^2\le q_z^a\) and \(|\iota\alpha^{-1}|^2\le q_z^{-a}\), and equality holds.

For a sheaf that is already lisse on the curve, the inertia is trivial, and Lemma \hyperref[bp6]{15.2} identifies the stalk of \(j_*\) with the original stalk. The first assertion, applied at each of the finitely many points outside the dense open, then gives the second.
\end{proof}

% text: prose-15-06
\phantomsection\label{bp18}%
The proof of Proposition \hyperref[bp13]{15.5}, together with Lemma \bpout{14.20}, also gives the extreme degrees of the cohomology of a middle extension on an affine curve.

% object: lemma-extreme-degrees-affine-curve (blueprint/chap-Pure-sheaves-on-curves/sec-Weights-of-middle-extensions/lemma-extreme-degrees-affine-curve/_section.md)
\phantomsection\label{bp19}%
\begin{bpstatementit}{Lemma~15.7}{ (The extreme degrees on an affine curve)}
Let \(j:U_0\hookrightarrow C_0\) be a dense open in a smooth geometrically connected affine curve, let \(F\) be lisse on \(U_0\), and put \(V=F_{\bar\eta}\). Then \(H_c^i(C,j_*F)=0\) for \(i\ne1,2\), and there is a canonical identification

\[
H_c^2(C,j_*F)=V_{\pi_1(U)}(-1).
\]

If \(F\) is pure of weight \(w\), the group in degree two, when nonzero, is pure of weight \(w+2\).
\end{bpstatementit}

% proof: proof/lemma-extreme-degrees-affine-curve
\begin{proof}[{Proof}]\phantomsection\label{bp20}%
The vanishing of \(H^0_c\) was shown in the proof of Proposition \hyperref[bp13]{15.5}. For the higher degrees, let \(i\) be the inclusion of \(C_0-U_0\). Since sheaves with finite support have compact-support cohomology only in degree zero, the localization sequence for \(0\to j_!F\to j_*F\to i_*i^*j_*F\to0\) identifies \(H^r_c(C,j_*F)\) with \(H_c^r(U,F)\) for \(r\ge2\). Lemma \bpout{14.20} gives the value and the weight of this group for \(r=2\), and its vanishing for \(r>2\).
\end{proof}

% text: prose-15-07
\phantomsection\label{bp21}%
For the perverse t-structure in Chapter \bpout{16} we need the same range of degrees for arbitrary constructible sheaves.

% object: lemma-cohomological-range-curve (blueprint/chap-Pure-sheaves-on-curves/sec-Weights-of-middle-extensions/lemma-cohomological-range-curve/_section.md)
\phantomsection\label{bp22}%
\begin{bpstatementit}{Lemma~15.8}{ (Cohomological range on a curve)}
Let \(C\) be a smooth curve over an algebraically closed field and let \(F\) be a constructible sheaf on \(C\) with coefficients in \(E\) or \(\overline{\mathbf Q}_\ell\). Then \(H^i_c(C,F)=0\) for \(i\notin\{0,1,2\}\).
\end{bpstatementit}

% proof: proof/lemma-cohomological-range-curve
\begin{proof}[{Proof}]\phantomsection\label{bp23}%
Choose a dense open \(j:U\hookrightarrow C\) on which \(F\) is lisse, and let \(i\) be the inclusion of the finite complement. The sequence \(0\to j_!j^*F\to F\to i_*i^*F\to0\) gives a long exact sequence in which \(H^r_c(C,i_*i^*F)\) vanishes for \(r\ne0\), because the support is finite. It therefore suffices to show \(H^r_c(U,j^*F)=0\) for \(r\notin\{0,1,2\}\). For \(r<0\) there is nothing to prove, and for \(r>2\), degreewise Poincaré duality (Corollary \bpout{10.13}, with \(d=1\)) shows that \(H^r_c(U,j^*F)\) is dual to \(H^{2-r}(U,(j^*F)^\vee(1))\), which is zero because \((j^*F)^\vee\) is a sheaf and \(2-r<0\).
\end{proof}

% text: prose-15-08
\phantomsection\label{bp24}%
A constructible Weil sheaf on a curve \(C_0\) over \(k_0\) gives a constructible sheaf on \(C=C_0\times_{k_0}k\), so the lemma applies to it as well.


% object: sec-Growth-of-squared-traces (blueprint/chap-Pure-sheaves-on-curves/sec-Growth-of-squared-traces/_section.md), content hidden
\subsection*{\texorpdfstring{15.2\quad Growth of squared traces}{15.2 Growth of squared traces}}
\phantomsection\label{bp25}%
\addcontentsline{toc}{subsection}{\texorpdfstring{15.2\quad Growth of squared traces}{15.2 Growth of squared traces}}
% text: prose-15-09
\phantomsection\label{bp26}%
In Chapter \bpout{17} we know the trace function of the Fourier transform \(M'\) exactly, but not its weight. For a pure sheaf the trace function determines the weight: it is the growth rate of the sums of squared traces (Theorem \hyperref[bp41]{15.14}).

% object: proposition-twisting-weight-zero-adding-dual (blueprint/chap-Pure-sheaves-on-curves/sec-Growth-of-squared-traces/proposition-twisting-weight-zero-adding-dual/_section.md)
\phantomsection\label{bp27}%
\begin{bpstatementit}{Proposition~15.9}{ (Twisting to weight zero and adding the dual)}
Let \(\mathcal F\ne0\) be lisse and pure of weight \(w\) on a nonempty finite-type scheme over \(k_0\).  There is \(b\in\overline{\mathbf Q}_\ell^\times\) with \(|\iota(b)|=q^{-w/2}\).  Put

\[
\mathcal F^0=\mathcal F\otimes\mathcal B_b,
\qquad \mathcal E=\mathcal F^0\oplus(\mathcal F^0)^\vee.
\]

Then \(\mathcal E\) is real and pure of weight zero, contains \(\mathcal F^0\) as a direct summand, and at points over \(k_n\) its trace is twice the real part of the trace of \(\mathcal F^0\).
\end{bpstatementit}

% proof: proof/proposition-twisting-weight-zero-adding-dual
\begin{proof}[{Proof}]\phantomsection\label{bp28}%
A twist to weight zero puts all local eigenvalues on the unit circle, where inversion is complex conjugation. Adding the dual therefore makes each multiset of eigenvalues stable under conjugation.

To find the twist, choose a closed point \(x\), of degree \(d\), at which the stalk is nonzero, and a Frobenius eigenvalue \(\alpha\) at \(x\).  Let \(b\in\overline{\mathbf Q}_\ell\) be a \(d\)th root of \(\alpha^{-1}\).  By purity, \(|\iota(b)|=|\iota(\alpha)|^{-1/d}=q^{-w/2}\), and Proposition \bpout{13.7} shows that \(\mathcal F^0\) is pure of weight zero.

Its local eigenvalues \(\beta_j\) therefore have absolute value one after applying \(\iota\). The local eigenvalues of its dual are the \(\beta_j^{-1}\), and \(\iota(\beta_j^{-1})=\overline{\iota(\beta_j)}\). So the union of the two multisets is stable under complex conjugation, and \(\mathcal E\) is real and pure. In the same way, at a point \(x\) over \(k_n\) the eigenvalues of \(\operatorname{Frob}_{x,n}\) on \((\mathcal F^0)^\vee\) become, after \(\iota\), the complex conjugates of those on \(\mathcal F^0\). Hence \(\iota\) of the trace of \((\mathcal F^0)^\vee\) is the complex conjugate of \(\iota\) of the trace of \(\mathcal F^0\), and the trace of \(\mathcal E\) is twice its real part.
\end{proof}

% text: prose-15-10
\phantomsection\label{bp29}%
The number \(\nu\) below is the growth exponent of the sums of squared traces, shifted by \(1\) so that it equals the weight. We take a \(\limsup\) because the sums can vanish for infinitely many \(n\), as the example of \(\mathcal B_b\oplus\mathcal B_{-b}\) at the start of this chapter shows.

% object: definition-growth-squared-traces (blueprint/chap-Pure-sheaves-on-curves/sec-Growth-of-squared-traces/definition-growth-squared-traces/_section.md)
\phantomsection\label{bp30}%
\begin{bpstatement}{Definition~15.10}{ (Growth of squared traces)}
Let \(K\) be a constructible Weil complex on an affine curve \(X_0\). For \(n\ge1\) and \(x\in X_0(k_n)\), let \(\operatorname{Frob}_{x,n}\) be the Frobenius acting on the stalk \(K_{\bar x}\) (Theorem \bpout{12.24}), and set

\[
t_{K,n}(x)=\sum_i(-1)^i\operatorname{Tr}
(\operatorname{Frob}_{x,n}\mid\mathcal H^i(K)_{\bar x}),\qquad
S_n(K)=\sum_{x\in X_0(k_n)}|\iota t_{K,n}(x)|^2.
\]

Put

\[
\nu(K)=\limsup_{n\to\infty}\frac{\log_q S_n(K)}n-1,
\qquad \log_q0=-\infty.
\]
\end{bpstatement}

% text: prose-15-11
\phantomsection\label{bp31}%
A cohomological shift changes the weight of a complex in the sense of Definition \bpout{13.1}. But it only multiplies the trace function by a sign, and so it leaves \(\nu\) unchanged.

In the theorem below we find \(\nu\) from the radius of convergence of a generating series, which is a logarithmic derivative of an \(L\)-function.

% object: lemma-radius-squared-traces (blueprint/chap-Pure-sheaves-on-curves/sec-Growth-of-squared-traces/lemma-radius-squared-traces/_section.md)
\phantomsection\label{bp32}%
\begin{bpstatementit}{Lemma~15.11}{ (Radius of convergence)}
Let \(K\) be a constructible Weil complex on an affine curve \(X_0\). The series \(\varphi_K(t)=\sum_{n\ge1}S_n(K)t^{n-1}\) has radius of convergence \(q^{-\nu(K)-1}\), where this number is read as \(+\infty\) if \(\nu(K)=-\infty\) and as \(0\) if \(\nu(K)=+\infty\).
\end{bpstatementit}

% proof: proof/lemma-radius-squared-traces
\begin{proof}[{Proof}]\phantomsection\label{bp33}%
By the Cauchy–Hadamard formula the radius is the inverse of \(\limsup_nS_n(K)^{1/n}=q^{\nu(K)+1}\); the shift of the index by one does not change it.
\end{proof}

% text: prose-15-12
\phantomsection\label{bp34}%
To compute \(\nu\) we also need upper bounds for lisse sheaves and for complexes with finite support, as well as the effect of twisting by \(\mathcal B_b\).

% object: lemma-bounds-squared-traces (blueprint/chap-Pure-sheaves-on-curves/sec-Growth-of-squared-traces/lemma-bounds-squared-traces/_section.md)
\phantomsection\label{bp35}%
\begin{bpstatementit}{Lemma~15.12}{ (Bounds for squared traces)}
Let \(X_0\) be a smooth affine curve over \(k_0\), and let \(a\) be a real number.

\begin{enumerate}
\item If \(F\) is a lisse Weil sheaf of rank \(r\) on \(X_0\), punctually of weights \(\le a\), then \(S_n(F)\le Dr^2q^{n(a+1)}\) for all \(n\), with \(D\) as in Lemma \hyperref[bp10]{15.4}, and hence \(\nu(F)\le a\).
\item If \(K\) is a constructible Weil complex on \(X_0\) supported at finitely many closed points, and every eigenvalue \(\alpha\) of \(\operatorname{Frob}_x\) on every \(\mathcal H^i(K)_{\bar x}\) satisfies \(|\iota(\alpha)|\le q_x^{a/2}\), then there is a constant \(C\) with \(S_n(K)\le Cq^{na}\) for all \(n\), and hence \(\nu(K)\le a-1\).
\item For every constructible Weil complex \(K\) on \(X_0\) and every \(b\in\overline{\mathbf Q}_\ell^\times\), we have \(\nu(K\otimes\mathcal B_b)=\nu(K)+\lambda_\iota(b)\), with \(\mathcal B_b\) as in Proposition \bpout{13.7}.
\end{enumerate}
\end{bpstatementit}

% proof: proof/lemma-bounds-squared-traces
\begin{proof}[{Proof}]\phantomsection\label{bp36}%
For (1), the trace at each point of \(X_0(k_n)\) has absolute value at most \(rq^{na/2}\), and Lemma \hyperref[bp10]{15.4} gives \(S_n\le Dr^2q^{n(a+1)}\). For (2), Lemma \hyperref[bp10]{15.4} bounds the number of points of the support with values in \(k_n\) independently of \(n\), so it suffices to bound the trace at each of them. Since the support consists of finitely many points with finite-dimensional stalks, the sum \(c\) of the ranks of the stalks of the \(\mathcal H^i(K)\) is finite, and at each such point the trace has absolute value at most \(cq^{na/2}\). For (3), the trace of \(\mathcal B_b\) at a \(k_n\)-point is \(b^n\), so tensoring by \(\mathcal B_b\) multiplies \(S_n\) by \(|\iota(b)|^{2n}=q^{n\lambda_\iota(b)}\).
\end{proof}

% text: prose-15-13
\phantomsection\label{bp37}%
A correction at finitely many points of small weight does not change \(\nu\). We use this to discard the boundary terms in Corollary \hyperref[bp44]{15.15}.

% object: lemma-growth-finite-corrections (blueprint/chap-Pure-sheaves-on-curves/sec-Growth-of-squared-traces/lemma-growth-finite-corrections/_section.md)
\phantomsection\label{bp38}%
\begin{bpstatementit}{Lemma~15.13}{ (Finite corrections)}
Let \(K\) and \(L\) be constructible Weil complexes on an affine curve \(X_0\), and let \(a\) be a real number. Suppose that there is a constant \(C\) with \(\sum_{x\in X_0(k_n)}|\iota(t_{K,n}-t_{L,n})(x)|^2\le Cq^{na}\) for all \(n\). If \(\max(\nu(K),\nu(L))>a-1\), then \(\nu(K)=\nu(L)\).
\end{bpstatementit}

% proof: proof/lemma-growth-finite-corrections
\begin{proof}[{Proof}]\phantomsection\label{bp39}%
The inequality \(|z+u|^2\le2|z|^2+2|u|^2\), applied to \(t_{L,n}\) and the difference \(t_{K,n}-t_{L,n}\), gives \(S_n(K)\le2S_n(L)+2Cq^{na}\), so that \(\nu(K)\le\max(\nu(L),a-1)\), and by symmetry \(\nu(L)\le\max(\nu(K),a-1)\).  If one of \(\nu(K)\), \(\nu(L)\) exceeds \(a-1\), these two inequalities force equality.
\end{proof}

% text: prose-15-14
\phantomsection\label{bp40}%
By counting points, Lemma lemma-bounds-squared-traces gives \(\nu\le w\) for a sheaf pure of weight \(w\). A priori, however, the squared traces could be much smaller than this trivial bound, through cancellation among the eigenvalues at each point. The next theorem shows that on average they are not.

% object: theorem-weights-sums-squared-traces (blueprint/chap-Pure-sheaves-on-curves/sec-Growth-of-squared-traces/theorem-weights-sums-squared-traces/_section.md)
\phantomsection\label{bp41}%
\begin{bpstatementit}{Theorem~15.14}{ (Weights and sums of squared traces)}
Let \(X_0/\mathbf F_q\) be a smooth geometrically connected affine curve. For every nonzero pure lisse Weil sheaf \(\mathcal G\) of weight \(w\),

\[
\nu(\mathcal G)=w.
\]
\end{bpstatementit}

% proof: proof/theorem-weights-sums-squared-traces
\begin{proof}[{Proof}]\phantomsection\label{bp42}%
Lemma lemma-bounds-squared-traces gives \(\nu(\mathcal G)\le w\), so, by Lemma \hyperref[bp32]{15.11}, we must show that the radius of \(\varphi_{\mathcal G}\) is not larger than \(q^{-w-1}\). For a real sheaf we find a pole of the logarithmic derivative of an \(L\)-function on that circle, and Proposition \hyperref[bp27]{15.9} reduces the general case to this one by twisting to weight zero and adding the dual.

\emph{The real case.} Suppose first that \(\mathcal G\) is real, and put \(\mathcal T=\mathcal G\otimes\mathcal G\). The squared traces of \(\mathcal G\) are then the traces of \(\mathcal T\), and we express their generating series through the \(L\)-function of \(\mathcal T\). Since \(\iota t_{\mathcal G,n}(x)\) is real, we have \(\iota t_{\mathcal T,n}(x)=|\iota t_{\mathcal G,n}(x)|^2\) for \(x\in X_0(k_n)\). The Euler product and the trace formula then give

\[
\varphi_{\mathcal G}(t)=\frac{L'(t)}{L(t)},\qquad
L(t)=\frac{\iota\det(1-Ft\mid H_c^1(X,\mathcal T))}
{\iota\det(1-Ft\mid H_c^2(X,\mathcal T))}.
\]

For the first equality, group the local logarithms by degree. Then the coefficient of \(t^{n-1}\) in \(L'/L\) is the sum of the traces over \(X_0(k_n)\), which is \(S_n(\mathcal G)\). As \(\mathcal T\) is pure of weight \(2w\), Lemma \bpout{14.20} shows that all roots of the denominator have absolute value \(R=q^{-w-1}\).

We show that the function \(L\) has a pole on the circle \(|t|=R\). This is not automatic, since a priori all the eigenvalues on \(H^2_c\) could cancel against eigenvalues on \(H^1_c\), but positivity rules this out. Indeed, by Lemma \bpout{14.22}, each local Euler factor of \(\mathcal T\) has nonnegative Taylor coefficients, so the coefficients of the full product dominate those of any one factor. Now choose a closed point of the curve. Since \(\mathcal G\) has positive rank and Frobenius is invertible, the Euler factor of \(\mathcal T\) at this point is the reciprocal of a nonconstant polynomial. It has infinitely many nonzero Taylor coefficients, because a polynomial cannot be inverse to a nonconstant polynomial. Hence the full Euler product is not a polynomial, and the rational function \(L\) has a pole, which by the bound on the denominator lies on \(|t|=R\).

A pole of \(L\) is also a pole of its logarithmic derivative: if \(L(t)=(t-t_0)^{-m}u(t)\) near \(t_0\), with \(m>0\) and \(u(t_0)\ne0\), then \(L'/L=-m/(t-t_0)+u'/u\). So the radius of convergence of \(\varphi_{\mathcal G}\) is exactly \(R\), and Lemma \hyperref[bp32]{15.11} gives \(\nu(\mathcal G)=w\).

\emph{The general case.} Let \(\mathcal G\) be an arbitrary pure sheaf, and let \(\mathcal G^0=\mathcal G\otimes\mathcal B_b\) and \(\mathcal E\) be as in Proposition \hyperref[bp27]{15.9}. The sheaf \(\mathcal G^0\) is pure of weight zero, and by Lemma lemma-bounds-squared-traces we have \(\nu(\mathcal G^0)=\nu(\mathcal G)+\lambda_\iota(b)=\nu(\mathcal G)-w\), so it suffices to show that \(\nu(\mathcal G^0)=0\). The sheaf \(\mathcal E=\mathcal G^0\oplus(\mathcal G^0)^\vee\) is nonzero, real and pure of weight zero, so by the case already treated \(\nu(\mathcal E)=0\).

This lower bound passes to \(\mathcal G^0\), because the trace of \(\mathcal E\) is twice the real part of that of \(\mathcal G^0\). In other words, for every \(n\) and \(x\in X_0(k_n)\) we have \(\iota t_{\mathcal E,n}(x)=2\operatorname{Re}(\iota t_{\mathcal G^0,n}(x))\). Hence \(S_n(\mathcal E)\le4S_n(\mathcal G^0)\), and so \(\nu(\mathcal G^0)\ge0\). Since the opposite inequality comes from Lemma lemma-bounds-squared-traces, we obtain \(\nu(\mathcal G^0)=0\), that is, \(\nu(\mathcal G)=w\).
\end{proof}

% text: prose-15-15
\phantomsection\label{bp43}%
Under the hypothesis of Proposition \bpout{16.17}, the Fourier transform of \((j_*M)[1]\) is again a middle extension. But Plancherel applies only to complexes on the whole line, boundary points included, so we need the theorem for \(j_*\) and \(j_!\) on the larger curve. This costs nothing: the boundary points do not change \(\nu\), because their weights are small.

% object: corollary-sums-squared-traces-j-shriek (blueprint/chap-Pure-sheaves-on-curves/sec-Growth-of-squared-traces/corollary-sums-squared-traces-j-shriek/_section.md)
\phantomsection\label{bp44}%
\begin{bpstatementit}{Corollary~15.15}{ (Sums of squared traces of $j_!$ and $j_*$)}
Let \(j:U_0\hookrightarrow X_0\) be a dense open of a smooth geometrically connected affine curve, and let \(\mathcal G\) be a nonzero pure lisse Weil sheaf of weight \(w\) on \(U_0\). Then

\[
\nu(j_!\mathcal G)=\nu(j_*\mathcal G)=w.
\]
\end{bpstatementit}

% proof: proof/corollary-sums-squared-traces-j-shriek
\begin{proof}[{Proof}]\phantomsection\label{bp45}%
For \(j_!\) the trace function is that of \(\mathcal G\) extended by zero, and for \(j_*\) it differs from this at finitely many points of small weight, which Lemma \hyperref[bp38]{15.13} allows us to ignore.

The open curve \(U_0\) is again smooth, geometrically connected and affine, and Theorem \hyperref[bp41]{15.14} on \(U_0\) gives \(\nu(\mathcal G)=w\). Extension by zero only adds zero traces, so \(\nu(j_!\mathcal G)=w\). The trace functions of \(j_*\) and of the extension by zero differ only at the points of \(X_0-U_0\), by the trace function of the boundary stalks. These stalks have weights at most \(w\) by Proposition \hyperref[bp13]{15.5}, so if \(i\) denotes the inclusion of \(X_0-U_0\), then Lemma lemma-bounds-squared-traces, applied to the sheaf \(i_*i^*j_*\mathcal G\), bounds the squared sum of the difference by \(Cq^{nw}\). Since \(\nu(j_!\mathcal G)=w>w-1\), Lemma \hyperref[bp38]{15.13} shows that \(\nu(j_*\mathcal G)\) is still \(w\).
\end{proof}



\clearpage\phantomsection\hypertarget{bp-outside}{}%
\section*{Outside this export}
\addcontentsline{toc}{section}{Outside this export}
References marked \ensuremath{{}^\ast} are to parts of the blueprint this document leaves out:
\begin{description}
\item[Corollary 10.13] Degreewise duality over fields, in Ⅱ Duality and the trace formula (\bpslug{corollary-{}\allowbreak{}degreewise-{}\allowbreak{}duality-{}\allowbreak{}over-{}\allowbreak{}fields})
\item[Theorem 12.24] Grothendieck–Lefschetz trace formula, in Ⅱ Duality and the trace formula (\bpslug{theorem-{}\allowbreak{}grothendieck-{}\allowbreak{}lefschetz-{}\allowbreak{}trace-{}\allowbreak{}formula})
\item[Definition 13.1] Weights, purity and mixedness, in Ⅲ Weights and the Riemann hypothesis (\bpslug{definition-{}\allowbreak{}weights})
\item[Example 13.4] The Kummer extension, in Ⅲ Weights and the Riemann hypothesis (\bpslug{example-{}\allowbreak{}kummer-{}\allowbreak{}extension})
\item[Proposition 13.7] Weights under the usual operations, in Ⅲ Weights and the Riemann hypothesis (\bpslug{proposition-{}\allowbreak{}weights-{}\allowbreak{}under-{}\allowbreak{}usual-{}\allowbreak{}operations})
\item[Lemma 14.20] Degree zero and top cohomology, in Ⅲ Weights and the Riemann hypothesis (\bpslug{lemma-{}\allowbreak{}degree-{}\allowbreak{}zero-{}\allowbreak{}top-{}\allowbreak{}cohomology})
\item[Definition 14.21] Real sheaves, in Ⅲ Weights and the Riemann hypothesis (\bpslug{definition-{}\allowbreak{}real-{}\allowbreak{}sheaves})
\item[Lemma 14.22] Positivity of Euler factors, in Ⅲ Weights and the Riemann hypothesis (\bpslug{lemma-{}\allowbreak{}positivity-{}\allowbreak{}euler-{}\allowbreak{}factors})
\item[Chapter 16] The Fourier transform, in Ⅲ Weights and the Riemann hypothesis (\bpslug{chap-{}\allowbreak{}The-{}\allowbreak{}Fourier-{}\allowbreak{}transform})
\item[Proposition 16.17] The Fourier transform of a middle extension, in Ⅲ Weights and the Riemann hypothesis (\bpslug{proposition-{}\allowbreak{}fourier-{}\allowbreak{}transform-{}\allowbreak{}middle-{}\allowbreak{}extension})
\item[Chapter 17] The Deligne–Laumon theorem for curves, in Ⅲ Weights and the Riemann hypothesis (\bpslug{chap-{}\allowbreak{}The-{}\allowbreak{}Deligne-{}\allowbreak{}Laumon-{}\allowbreak{}theorem-{}\allowbreak{}for-{}\allowbreak{}curves})
\end{description}
\clearpage\phantomsection%
\section*{References}
\addcontentsline{toc}{section}{References}
\bpbibitem{BBD82}{bp46}{A. Beilinson, J. Bernstein, and P. Deligne. \emph{Faisceaux pervers}, volume 100 of \emph{Astérisque}. 1982.}

\bpbibitem{BS15}{bp47}{B. Bhatt and P. Scholze. The pro-étale topology for schemes. \emph{Astérisque}, 369:99–201, 2015. \href{https://people.mpim-bonn.mpg.de/scholze/proetale.pdf}{Author's manuscript}.}

\bpbibitem{Con07}{bp48}{B. Conrad. Deligne's notes on Nagata compactifications. \emph{Journal of the Ramanujan Mathematical Society}, 22:205–257, 2007. \href{https://math.stanford.edu/\~conrad/papers/nagatafinal.pdf}{Author's manuscript}.}

\bpbibitem{Del74}{bp49}{P. Deligne. La conjecture de Weil. I. \emph{Publ. Math. IHÉS}, 43:273–307, 1974.}

\bpbibitem{Del80}{bp50}{P. Deligne. La conjecture de Weil. II. \emph{Publ. Math. IHÉS}, 52:137–252, 1980.}

\bpbibitem{Den22}{bp51}{C. Deninger. There is no “Weil-”cohomology theory with real coefficients for arithmetic curves. \emph{Ann. Sc. Norm. Super. Pisa Cl. Sci.}, 2022. To appear. arXiv:\href{https://arxiv.org/abs/2204.02714}{2204.02714}.}

\bpbibitem{Deu41}{bp52}{M. Deuring. Die Typen der Multiplikatorenringe elliptischer Funktionenkörper. \emph{Abh. Math. Sem. Univ. Hamburg}, 14:197–272, 1941.}

\bpbibitem{Eke90}{bp53}{T. Ekedahl. On the adic formalism. In \emph{The Grothendieck Festschrift, Vol. II}, volume 87 of \emph{Progress in Mathematics}, pages 197–218. Birkhäuser, 1990.}

\bpbibitem{Has36}{bp54}{H. Hasse. Zur Theorie der abstrakten elliptischen Funktionenkörper. I, II, III. \emph{J. Reine Angew. Math.}, 175:55–62, 69–88, 193–208, 1936.}

\bpbibitem{Hum72}{bp55}{J. E. Humphreys. \emph{Introduction to Lie Algebras and Representation Theory}, volume 9 of \emph{Graduate Texts in Mathematics}. Springer, 1972.}

\bpbibitem{Ill94}{bp56}{L. Illusie. Autour du théorème de monodromie locale. \emph{Astérisque}, 223:9–57, 1994. §4. \href{https://www.numdam.org/item/AST_1994__223__9_0/}{Original article}.}

\bpbibitem{Kat88}{bp57}{N. M. Katz. Travaux de Laumon. \emph{Astérisque}, 161–162:105–132, 1988. Séminaire Bourbaki, exp. 691.}

\bpbibitem{KW01}{bp58}{R. Kiehl and R. Weissauer. \emph{Weil conjectures, perverse sheaves and \(\ell\)-adic Fourier transform}, volume 42 of \emph{Ergebnisse der Mathematik}. Springer, 2001.}

\bpbibitem{Lau87}{bp59}{G. Laumon. Transformation de Fourier, constantes d'équations fonctionnelles et conjecture de Weil. \emph{Publ. Math. IHÉS}, 65:131–210, 1987.}

\bpbibitem{Mil86a}{bp60}{J. S. Milne. \href{https://www.jmilne.org/math/xnotes/AVs.pdf}{Abelian Varieties}. In \emph{Arithmetic Geometry}. 1986. Corrected online version.}

\bpbibitem{Mil86b}{bp61}{J. S. Milne. \href{https://www.jmilne.org/math/xnotes/JVs.pdf}{Jacobian Varieties}. In \emph{Arithmetic Geometry}. 1986. Corrected online version.}

\bpbibitem{Ran39}{bp62}{R. A. Rankin. Contributions to the theory of Ramanujan's function \(\tau(n)\) and similar arithmetical functions. II. \emph{Proc. Cambridge Philos. Soc.}, 35:357–372, 1939.}

\bpbibitem{SGA4}{bp63}{M. Artin, A. Grothendieck, and J.-L. Verdier, editors. \emph{Théorie des topos et cohomologie étale des schémas}, volumes 269, 270, 305 of \emph{Lecture Notes in Mathematics}. Springer, 1972–1973. SGA 4.}

\bpbibitem{SGA4½}{bp64}{P. Deligne. \emph{Cohomologie étale}, volume 569 of \emph{Lecture Notes in Mathematics}. Springer, 1977. SGA 4½; \emph{Théorèmes de finitude}, §3, and \emph{Rapport sur la formule des traces}, §3, and \emph{Sommes trigonométriques}, §1.}

\bpbibitem{SGA5}{bp65}{A. Grothendieck, editor. \emph{Cohomologie \(\ell\)-adique et fonctions \(L\)}, volume 589 of \emph{Lecture Notes in Mathematics}. Springer, 1977. SGA 5.}

\bpbibitem{SGA7II}{bp66}{P. Deligne and N. Katz, editors. \emph{Groupes de monodromie en géométrie algébrique}, volume 340 of \emph{Lecture Notes in Mathematics}. Springer, 1973. SGA 7 II; Exposé XIII, (2.1.7).}

\bpbibitem{Stacks}{bp67}{The Stacks Project Authors. \href{https://stacks.math.columbia.edu}{The Stacks Project}. Individual results are cited by their stable tags in the text.}

\bpbibitem{Var05}{bp68}{Y. Varshavsky. Lefschetz–Verdier trace formula and a generalization of a theorem of Fujiwara. 2005. Version 2. arXiv:\href{https://arxiv.org/abs/math/0505564}{math/0505564}. \href{https://arxiv.org/abs/math/0505564v2}{Author's preprint}.}

\bpbibitem{Wei48a}{bp69}{A. Weil. \emph{Sur les courbes algébriques et les variétés qui s'en déduisent}. Hermann, Paris, 1948.}

\bpbibitem{Wei48b}{bp70}{A. Weil. \emph{Variétés abéliennes et courbes algébriques}. Hermann, Paris, 1948.}

\bpbibitem{Wei49}{bp71}{A. Weil. Numbers of solutions of equations in finite fields. \emph{Bull. Amer. Math. Soc.}, 55:497–508, 1949.}

\end{document}
